% Relatorio da tarefa da Tabela III. NAO faz parte do build do repositorio:
% sai por e-mail. Todo numero vem de tabelas/gerado.tex, escrito por gerar.py.
% Os trechos de codigo sao recortes de lab.py e etrych.py: as linhas de codigo
% sao verbatim, trechos omitidos levam reticencias, e os comentarios foram
% traduzidos e encurtados para caber na pagina.
\documentclass[11pt,a4paper]{article}
\usepackage[T1]{fontenc}
\usepackage[utf8]{inputenc}
\usepackage{amsmath,amssymb}
\usepackage{booktabs}
\usepackage{graphicx}
\usepackage{xcolor}
\usepackage{listings}
\usepackage[margin=2.4cm]{geometry}
\usepackage[font=small,labelfont=bf,skip=4pt]{caption}
\usepackage[hidelinks]{hyperref}
% Keeps a table from drifting past the section heading that follows it: the
% a_- calibration table was floating to AFTER the positive-side subsection,
% so the subsection read before the table it builds on.
\usepackage[section]{placeins}

\definecolor{fundocodigo}{gray}{0.97}
\definecolor{comentario}{rgb}{0.30,0.45,0.30}
\definecolor{palavra}{rgb}{0.10,0.25,0.55}
\definecolor{texto}{rgb}{0.55,0.25,0.15}
\definecolor{regua}{gray}{0.75}

\lstdefinestyle{python}{
  language=Python,
  basicstyle=\ttfamily\scriptsize,
  keywordstyle=\color{palavra}\bfseries,
  commentstyle=\color{comentario}\itshape,
  stringstyle=\color{texto},
  backgroundcolor=\color{fundocodigo},
  frame=leftline, framesep=7pt, rulecolor=\color{regua},
  showstringspaces=false, breaklines=true, breakatwhitespace=true,
  columns=fullflexible, keepspaces=true, aboveskip=0.9em, belowskip=0.7em,
  literate={á}{{\'a}}1 {ã}{{\~a}}1 {é}{{\'e}}1 {ê}{{\^e}}1 {í}{{\'i}}1
           {ó}{{\'o}}1 {õ}{{\~o}}1 {ú}{{\'u}}1 {ç}{{\c c}}1,
}
\lstset{style=python}

\newcommand{\cod}[1]{\texttt{\small #1}}

\renewcommand{\abstractname}{Resumo}
\renewcommand{\tablename}{Tabela}
\renewcommand{\figurename}{Figura}
\renewcommand{\refname}{Referências}
\renewcommand{\lstlistingname}{Código}
\newcommand{\hoje}{\number\day\ de \ifcase\month\or janeiro\or fevereiro\or
  março\or abril\or maio\or junho\or julho\or agosto\or setembro\or outubro\or
  novembro\or dezembro\fi\ de \number\year}

\newcommand{\FormaWell}{poço quadrado}
\newcommand{\FormaMpt}{P\"oschl-Teller}
\newcommand{\FormaGauss}{gaussiana}
\newcommand{\FormaLj}{Lennard-Jones}
\newcommand{\FormaCseis}{por $C_6$ sozinho}

%% generated by gerar.py -- do not edit

\newcommand{\TabEtrychIIICorpo}{%
$|1,1\rangle$ & 402.74 & 2.8 & 136 & $-830$ & $\pm$ & 40 & 900 & $\pm$ & 60 & $-0.1639$ \\
$|1,0\rangle$ & 472.33 & 2.8 & 137 & $-840$ & $\pm$ & 30 & 760 & $\pm$ & 70 & $-0.1631$ \\
$|1,-1\rangle$ & 33.5820 & 2.6 & 135 & $-908$ & $\pm$ & 11 & 876 & $\pm$ & 28 & $-0.1487$ \\
$|1,-1\rangle$ & 162.36 & 1.1 & 59 & $-780$ & $\pm$ & 70 & \multicolumn{3}{c}{---} & $-0.0756$ \\
$|1,-1\rangle$ & 561.14 & 2.5 & 132 & $-810$ & $\pm$ & 30 & 800 & $\pm$ & 100 & $-0.1630$ \\
}

\newcommand{\TabEtrychUnitarioCorpo}{%
\FormaWell & 1.233700550 & 1.000000000 & \ensuremath{-2.17992 \times 10^{-16}} & 1.000000000 \\
\FormaMpt & 1.000000000 & 2.000000000 & \ensuremath{-1.25557 \times 10^{-13}} & 1.000000000 \\
\FormaGauss & 1.342002325 & 1.435224423 & \ensuremath{3.01187 \times 10^{-14}} & 1.000000000 \\
\FormaLj & 0.264624012 & 0.000340682 & \ensuremath{1.40028 \times 10^{-15}} & 1.000000000 \\
}

\newcommand{\TabEtrychAminusCorpo}{%
136 & $-830$ & \FormaWell & 1.09639 & 0.00781181 & $-830.000$ & 136.000 & 0 \\
 &  & \FormaMpt & 0.897532 & 0.0160072 & $-830.000$ & 136.000 & 0 \\
 &  & \FormaGauss & 1.1985 & 0.011357 & $-830.000$ & 136.000 & 0 \\
 &  & \FormaLj & \ensuremath{6.03172 \times 10^{7}} & \ensuremath{2.87324 \times 10^{17}} & $-830.000$ & 136.000 & 0 \\
\addlinespace
137 & $-840$ & \FormaWell & 1.09693 & 0.0077528 & $-840.000$ & 137.000 & 0 \\
 &  & \FormaMpt & 0.897931 & 0.0158849 & $-840.000$ & 137.000 & 0 \\
 &  & \FormaGauss & 1.19906 & 0.0112706 & $-840.000$ & 137.000 & 0 \\
 &  & \FormaLj & \ensuremath{6.22169 \times 10^{7}} & \ensuremath{3.10394 \times 10^{17}} & $-840.000$ & 137.000 & 0 \\
\addlinespace
135 & $-908$ & \FormaWell & 1.10724 & 0.00782909 & $-908.000$ & 135.000 & 0 \\
 &  & \FormaMpt & 0.905607 & 0.0160129 & $-908.000$ & 135.000 & 0 \\
 &  & \FormaGauss & 1.20983 & 0.0113709 & $-908.000$ & 135.000 & 0 \\
 &  & \FormaLj & \ensuremath{6.06086 \times 10^{7}} & \ensuremath{2.89089 \times 10^{17}} & $-908.000$ & 135.000 & 0 \\
\addlinespace
59 & $-780$ & \FormaWell & 1.16431 & 0.0174533 & $-780.000$ & 59.000 & 0 \\
 &  & \FormaMpt & 0.948119 & 0.0353451 & $-780.000$ & 59.000 & 0 \\
 &  & \FormaGauss & 1.26942 & 0.0252155 & $-780.000$ & 59.000 & 0 \\
 &  & \FormaLj & \ensuremath{2.63133 \times 10^{6}} & \ensuremath{1.10114 \times 10^{14}} & $-780.000$ & 59.000 & 0 \\
\addlinespace
132 & $-810$ & \FormaWell & 1.09702 & 0.0080461 & $-810.000$ & 132.000 & 0 \\
 &  & \FormaMpt & 0.898001 & 0.0164856 & $-810.000$ & 132.000 & 0 \\
 &  & \FormaGauss & 1.19916 & 0.0116969 & $-810.000$ & 132.000 & 0 \\
 &  & \FormaLj & \ensuremath{5.36353 \times 10^{7}} & \ensuremath{2.14165 \times 10^{17}} & $-810.000$ & 132.000 & 0 \\
\addlinespace
}

\newcommand{\TabEtrychFamiliasCorpo}{%
\FormaWell & 0.877821 & 1.121244 & 27.7\% & 1.00000000 & 1.00000000 & 64.1\% \\
\FormaMpt & 0.822880 & 1.176102 & 42.9\% & 1.00000000 & 1.00000000 & 57.2\% \\
\FormaGauss & 0.847802 & 1.149652 & 35.6\% & 1.00000000 & 1.00000000 & 60.5\% \\
\FormaLj & 0.833849 & 1.193809 & 43.2\% & 0.99999999 & 1.00000006 & 22.9\% \\
\addlinespace
\quad {\itshape \FormaCseis} & 0.826206 & 1.206900 & 46.1\% & \multicolumn{2}{c}{---} & 13.1\% \\
}
\newcommand{\EtrychVunit}{1.2337005501}
\newcommand{\EtrychVexato}{1.2337005501}
\newcommand{\EtrychVerro}{\ensuremath{6.88338 \times 10^{-15}}}
\newcommand{\EtrychUnitInv}{\ensuremath{1.97842 \times 10^{-15}}}
\newcommand{\EtrychUnitRzero}{\ensuremath{4.68281 \times 10^{-11}}}
\newcommand{\EtrychAlvoA}{\ensuremath{2.88303 \times 10^{-12}}}
\newcommand{\EtrychAlvoRzero}{\ensuremath{1.01692 \times 10^{-8}}}
\newcommand{\EtrychAlvoBaixo}{-0.30}
\newcommand{\EtrychAlvoAlto}{0.30}
\newcommand{\EtrychLequeBorda}{6.5\%}
\newcommand{\EtrychLequeAmenos}{3.3\%}
\newcommand{\EtrychLequeRaso}{1.2\%}
\newcommand{\EtrychMinimoCseis}{2.1\%}
\newcommand{\EtrychDerivaWell}{27.7\%}
\newcommand{\EtrychForcaWell}{64.1\%}
\newcommand{\EtrychDerivaMpt}{42.9\%}
\newcommand{\EtrychForcaMpt}{57.2\%}
\newcommand{\EtrychDerivaGauss}{35.6\%}
\newcommand{\EtrychForcaGauss}{60.5\%}
\newcommand{\EtrychDerivaLj}{43.2\%}
\newcommand{\EtrychForcaLj}{22.9\%}
\newcommand{\EtrychDerivaCseis}{46.1\%}
\newcommand{\EtrychForcaCseis}{13.1\%}


\title{\vspace{-1.3cm}\textbf{Calibração de potenciais de dois corpos\\
sobre a Tabela III de Etrych \textit{et al.}}}
\author{Pedro Henrique Gesualdo Modesto\\[2pt]
\small Instituto de Física de São Carlos --- Universidade de São Paulo\\
\small Orientador: Lucas Madeira \quad Coorientadora: Patrícia C. M. Castilho}
\date{\small\hoje}

\begin{document}
\maketitle
\vspace{-0.5cm}
\begin{abstract}
\noindent
Para cada uma das cinco ressonâncias de Feshbach da Tabela III de
Etrych \textit{et al.} [1] são determinados os parâmetros de quatro formas de
potencial de dois corpos que reproduzem simultaneamente o alcance efetivo
$r_\mathrm{eff}$ e o comprimento de espalhamento $a_-$ da feição de Efimov: duas
equações a duas incógnitas por forma. Em seguida é construída a família que
leva de $a_-$ através da ressonância até o lado positivo, em duas prescrições.
A premissa de que tudo isso depende --- que uma colisão lenta guarda $a$ e
$r_0$ e mais nada do potencial --- é \emph{medida} aqui, não suposta: as quatro
formas, casadas nos dois números, concordam no deslocamento de fase em energia
finita a \EtrychPremissaMeio{} em $k r_0=\EtrychPremissaKmeio{}$, e o primeiro
termo que os dois números deixam de fora segue potência
\EtrychPremissaPotencia{} com coeficiente \EtrychPremissaCoef{} --- é o
parâmetro de forma, e não erro numérico. A energia do dímero concorda a
\EtrychDimeroLeque{}. A discordância que sobra entre as formas no mesmo ponto
--- \EtrychLequeAmenos{} em $r_0$ na faixa onde os $a_-$ se encontram --- mede
a dependência de forma que um cálculo de três corpos herdaria da escolha do
potencial, e é pequena comparada com o fator $1{,}3$ a $1{,}4$ que separa o
$a_-$ medido da previsão de van der Waals. Registra-se que $a_-$ é observável
de três corpos e entra como dado: constrói-se a \emph{entrada} de um cálculo de
três corpos, não uma previsão de $a_-$. A mesma calibração é feita na feição do
lado positivo, $a_p$, que o experimento reporta como a mais acessível das duas,
e lá a dependência de forma é igualmente pequena. Um resultado colateral
organiza o passo seguinte: $(a,r_0)$ não determina o potencial, e os ramos que
carregam dímeros profundos são distinguíveis pelo parâmetro de largura
$\eta_-$ que a mesma tabela publica.
\end{abstract}

\section{Objetivo}

Em energia baixa a colisão de dois átomos depende do potencial apenas por dois
comprimentos, $a$ e $r_0$. Em três corpos, o mesmo limite gera a torre de
Efimov, com invariância de escala discreta de razão $e^{\pi/s_0}\approx22{,}7$
($s_0=1{,}00624$): a simetria fixa as razões entre feições observáveis, mas não
suas posições absolutas, que dependem de uma condição de curta distância. A
universalidade de van der Waals prevê
$a_-\simeq-9{,}7\,r_\mathrm{vdW}\approx-630\,a_0$ para o potássio-39;
Etrych \textit{et al.} mediram $-13(1)\,r_\mathrm{vdW}$ nas quatro ressonâncias
que caracterizaram, e a Tabela III reúne essas quatro com uma quinta vinda de
outras referências.

Este trabalho faz o passo de dois corpos que antecede qualquer explicação dessa
discrepância: construir potenciais que sentam sobre os números medidos, para
que um cálculo de três corpos os receba com o alcance finito já embutido, em
vez de tratá-lo como correção perturbativa sobre alcance zero. São dois itens:
(i) calibrar, para cada linha e cada forma, os parâmetros que dão $a=a_-$ com
$r_0=r_\mathrm{eff}$, e os que dão unitariedade com $r_0=r_\mathrm{eff}$;
(ii) construir o caminho contínuo de $a_-$ até o lado positivo e medir o que
acontece com $r_0$ ao longo dele.

\section{Método}

Em onda $s$, com $\psi=u(r)/r$, a equação radial é
$-\frac{\hbar^2}{2\mu}u''+Vu=Eu$ com $u(0)=0$. O código divide energias e
potenciais por $\hbar^2/\mu$, de modo que $[V]=[E]=L^{-2}$ e a equação
integrada é $u''=2(V-E)u$; em energia zero, $u''=2Vu$. Fora do alcance do
potencial $V=0$ e $u$ é uma reta, de onde saem os dois observáveis:
\begin{equation}
  u(r)=C\left(1-\frac ra\right)
  \ \Longrightarrow\
  \frac1a=\frac{u'(R)}{u'(R)R-u(R)},
  \qquad
  r_0=2\int_0^\infty\!\left[\left(1-\frac ra\right)^{2}-u_0^{2}\right]dr .
  \label{eq:ar0}
\end{equation}
A convenção de sinal é a da Eq.~(A3) de [1], de modo que o $r_\mathrm{eff}$
tabulado \emph{é} o $r_0$ calculado aqui, sem conversão. Guarde a segunda
integral: o integrando é a diferença entre duas quantidades que se igualam
longe do potencial, e é ela que dita a arquitetura da Seção~\ref{sec:impl}.

Nessas unidades a equação não tem escala própria: sob $r\to\lambda r$ valem
$a\to\lambda a$ e $r_0\to\lambda r_0$, logo \textbf{a forma do potencial decide
apenas o adimensional $r_0/a$}. Cada linha da Tabela~\ref{tab:iii} é então uma
multiplicação de uma solução calculada uma única vez por forma, e as famílias
do item (ii) são universais: uma família em $r_0=1$ serve para as cinco linhas.

\section{Implementação}
\label{sec:impl}

O código está em \cod{lab.py} (dado $V(r)$, quanto valem $a$, $r_0$ e o número
de estados), \cod{etrych.py} (dados $r_\mathrm{eff}$ e $a_-$, quais parâmetros os
produzem), \cod{gerar.py} (as tabelas, as figuras e os dados deste relatório) e
\cod{test\_lab.py} (quarenta e sete conferências). Nas listagens abaixo as linhas de
código são verbatim; trechos omitidos levam reticências e os comentários foram
traduzidos.

\paragraph{O integrador.} Um só produz todos os números: Dormand--Prince de
oitava ordem, passo adaptativo, tolerância relativa $10^{-12}$. A decisão que
importa está na terceira componente do estado.

\begin{lstlisting}[caption={\cod{lab.py} --- o estado é $y=(u,u',\int u^2dr)$.
Carregar a integral da Eq.~\eqref{eq:ar0} \emph{dentro} do sistema faz com que
ela seja produzida pelo mesmo controle de erro que produz a função de onda, em
vez de por quadratura aplicada por cima: não resta malha para convergir. Os nós
saem por detecção de evento, e a integração é feita por trechos porque um passo
adaptativo que atravessa a descontinuidade do poço quadrado colapsa.}]
def _adaptativo(pot, E, rtol=RTOL, atol=ATOL, eventos=False):
    r_start = pot.r_min if pot.r_min > 0.0 else pot.R * R_START_FRACTION
    y0 = [0.0, 1.0, 0.0] if pot.r_min > 0.0 else [r_start, 1.0, 0.0]

    def campo(r, y):
        return (y[1], 2.0 * (float(pot.V(r)) - E) * y[0], y[0] * y[0])
    ...                                  # evento: zero de u, raiz refinada
    y = list(y0)
    for ini, fim in _trechos(pot, r_start, pot.R):
        sol = solve_ivp(campo, (ini, fim), y, method="DOP853",
                        rtol=rtol, atol=atol, **kwargs)
        y = list(sol.y[:, -1])
        ...
    return y, np.asarray(nos), r_start
\end{lstlisting}

\paragraph{A extração.} No raio de corte, valor e derivada dão os três
observáveis. A parte da reta da Eq.~\eqref{eq:ar0} é integrada analiticamente
sobre todo o intervalo e a parte de $u^2$ veio de dentro do sistema: nenhuma
quadratura sobrevive.

\begin{lstlisting}[caption={\cod{lab.py} --- perseguir $1/a$ e não $a$ percorre
todo o código: $a$ tem polos na ressonância, $1/a$ passa por zero.}]
def _uma_passada(pot):
    y_R, zeros, r_start = _adaptativo(pot, 0.0, eventos=True)
    u_R, slope, integral = y_R[0], y_R[1], y_R[2]

    inv_a = slope / (slope * pot.R - u_R)          # sem ramo, sem sinal a inventar
    a = 1.0 / inv_a

    C = (1.0 - pot.R / a) / u_R
    reta = pot.R - pot.R ** 2 / a + pot.R ** 3 / (3.0 * a ** 2)
    r0 = 2.0 * (reta - C * C * integral)

    nodes = int(np.sum(zeros > r_start * (1.0 + 1e-12)))
    if pot.R < a and abs(r0 / a) > UNITARITY:       # a reta contribui com um no
        nodes += 1                                  # em r = a, ate' a unitariedade
    return a, r0, nodes
\end{lstlisting}

A alternativa natural --- integrar a diferença $v=u-(1-r/a)$, que removeria o
cancelamento --- foi implementada e descartada: $v$ obedece à equação homogênea
no caroço, de modo que integrá-la para dentro através de uma parede repulsiva
seleciona a solução exponencialmente crescente. Medido no potencial HFD-B, $v$
alcançou $10^{240}$.

\paragraph{A cauda.} Truncar em $R$ custa pouco a $a$ (erro $\sim R^{-3}$) mas
custa a $r_0$ um termo que cai apenas como $1/R$ quando a cauda é de van der
Waals: cada década de corte compra um dígito. Empurrar $R$ também não resolve
--- passado $R\sim10^6$ a parte da reta carrega o arredondamento acumulado e
$r_0$ colapsa, medido. O corte não é escolhido: é eliminado.

\begin{lstlisting}[caption={\cod{lab.py} --- com $r_0(R)=r_0-K/R$, dois cortes a
um fator $F$ de distância eliminam o termo; um terceiro em $F^2R$ confere. Onde
a cauda é exponencial ou ausente os três já concordam e a extrapolação devolve
o mesmo número.}]
def _extrapolado(pot, fator=3.0):
    a1, r0_1, nos = _num_corte(pot, pot.R)
    a2, r0_2, nos2 = _num_corte(pot, fator * pot.R)
    r0 = (fator * r0_2 - r0_1) / (fator - 1.0)
    return a2, r0, nos2
\end{lstlisting}

A troca dessa política corrigiu um sistemático real: no corte fixo anterior os
alcances efetivos do Lennard-Jones saíam baixos em até $1{,}6\times10^{-3}$
relativo, muito acima do ruído do integrador.

\paragraph{O problema inverso.} A calibração é duas equações a duas incógnitas,
resolvida como duas buscas encaixadas --- a externa move a escala até $r_0$
fechar, a interna move a intensidade até $1/a$ fechar.

\begin{lstlisting}[caption={\cod{lab.py} --- o limiar é \emph{relativo}: a mesma
rotina é chamada para $r_0=1$ e para $r_0=136\,a_0$, e um limiar absoluto seria
duas exigências diferentes. \cod{nodes\_target} não é opcional: duas soluções
podem compartilhar $(a,r_0)$ e diferir em número de nós, e então não são o
mesmo estado físico.}]
def tune(name, a_target, r0_target, strength, scale, nodes_target=None):
    if math.isinf(a_target):
        inv_a_target = 0.0
    else:
        inv_a_target = 1.0 / a_target
    for step in range(25):
        strength = match_a(name, scale, strength, inv_a_target)
        a, r0, nodes = scattering(POTENTIALS[name](strength, scale))
        if abs(r0 / r0_target - 1.0) < 1e-8:
            break
        def r0_error(trial_scale):
            s = match_a(name, trial_scale, strength, inv_a_target)
            a2, r02, n2 = scattering(POTENTIALS[name](s, trial_scale))
            return r02 - r0_target
        scale = find_root(r0_error, scale)
    ...
    ok = abs(r0 / r0_target - 1.0) < 1e-7
    if nodes_target is not None and nodes != nodes_target:
        ok = False
    return strength, scale, a, r0, nodes, ok
\end{lstlisting}

Usa-se bisseção encaixada e não Newton em duas variáveis porque as curvas de
nível de $1/a$ e de $r_0$ no plano (intensidade, escala) são quase paralelas em
boa parte do domínio: o jacobiano seria mal condicionado exatamente ali. O
preço é de $646$ a $785$ chamadas do integrador por ponto, medidas nas quatro
formas; optou-se pela robustez. O limiar de $10^{-7}$ é onde o problema direto
deixa de ser reprodutível: medido, o ajuste do Lennard-Jones estaciona em
$4{,}0\times10^{-8}$, que é o ruído da extrapolação da cauda.

\paragraph{A escala, usada de propósito.} \cod{rescale} implementa
$r\to\lambda r$ com os expoentes ditados pela dimensão de cada coeficiente
($[C_6]=L^4$, $[C_{12}]=L^{10}$); \cod{unit\_solution} busca a solução unitária
uma vez por forma e guarda; \cod{unitary} e \cod{calibrate} só multiplicam.

\begin{lstlisting}[caption={\cod{etrych.py} --- a economia não é cosmética.
Ajustado direto em $r_\mathrm{eff}=136\,a_0$ o Lennard-Jones leva $144$~s e
\emph{falha} em três das cinco linhas, porque o chute inicial carrega $L^{10}$;
ajustado em $r_0=1$ e reescalado leva $1{,}4$~s.}]
def rescale(name, p1, p2, L):
    if name == "lj":
        return p1 * L ** 4, p2 * L ** 10
    return p1, p2 / L

def calibrate(name, row):          # a = a_-, r0 = r_eff, para uma linha
    p1, p2 = at_ratio(name, ratio(row))
    return rescale(name, p1, p2, row["r_eff"])
\end{lstlisting}

Ao caminhar de um ponto a outro da família, \cod{at\_ratio} parte da razão
vizinha já resolvida em vez da unitariedade, com o passo limitado: a solução se
move suavemente ao longo da família, e o vizinho é um chute muito melhor.
Medido sobre uma família de 25 pontos, $3{,}6$ vezes menos chamadas do
integrador, com parâmetros idênticos em $6\times10^{-14}$.

\paragraph{As duas famílias.} \cod{family\_retuned} reajusta os dois parâmetros
a cada ponto, travando $r_0$ por construção; \cod{family\_fixed\_range} congela
o alcance na solução unitária e varre só a intensidade, e aí $r_0$ é
\emph{resultado} --- o quanto ele anda é a medida da dependência de forma. As
duas são lidas na \emph{mesma} lista de $1/a$: varrer a profundidade
diretamente seria a comparação errada, porque o mesmo passo em profundidade é
um passo diferente em $1/a$ para cada forma.

Para o Lennard-Jones, ``aumentar a intensidade'' admite duas leituras. Com
$V=\tfrac12(C_{12}/r^{12}-C_6/r^6)$, o mínimo fica em $r_m^6=2C_{12}/C_6$ com
profundidade $C_6^2/8C_{12}$: multiplicar \emph{ambos} os coeficientes pelo
mesmo fator deixa $r_m$ onde estava e multiplica a profundidade, que é o que as
outras três formas fazem; mexer só em $C_6$ aprofunda \emph{e} desloca o
mínimo. \cod{sweep\_c6\_alone} existe para medir a diferença em vez de
afirmá-la.

\section{Resultados}
\label{sec:resultados}

\begin{table}[htbp]
\centering
\caption{Tabela III de [1], transcrita. Comprimentos em raios de Bohr. Três
ressalvas da fonte: o $r_\mathrm{eff}$ não é medido, vem do cálculo de canais
acoplados dos autores perto de $B_\mathrm{res}$; a linha de $33{,}5820$~G é
reproduzida ali de outras referências; e $a_p$ é o símbolo do artigo e
\emph{não} é o $a_+$ dele --- a modulação log-periódica do lado positivo tem
mínimos em $a_+e^{j\pi/s_0}$ e máximos em $a_pe^{j\pi/s_0}$, com
$a_p=a_+e^{\pi/(2s_0)}$. A razão universal testada é $a_p/a_-=-1$.}
\label{tab:iii}
\small
\begin{tabular}{lrrrr@{\,}c@{\,}lr@{\,}c@{\,}lr}
\toprule
estado & $B_\mathrm{res}$ (G) & $s_\mathrm{res}$ & $r_\mathrm{eff}$ &
\multicolumn{3}{c}{$a_-$} & \multicolumn{3}{c}{$a_p$} & $r_\mathrm{eff}/a_-$ \\
\midrule
\TabEtrychIIICorpo
\bottomrule
\end{tabular}
\end{table}

\begin{table}[htbp]
\centering
\caption{\textbf{Item (i).} Acima: as quatro soluções unitárias em $r_0=1$,
calculadas uma vez por forma; toda a calibração unitária das cinco linhas sai
delas por multiplicação. Abaixo: a calibração em $a=a_-$ com $r_0=r_\mathrm{eff}$,
vinte linhas, em raios de Bohr. Em todos os casos o ramo é o mais raso, sem
estados ligados profundos --- escolha discutida na Seção~\ref{sec:ramos}.}
\label{tab:calib}
\small
\begin{tabular}{lrrrr}
\toprule
forma & primeiro parâmetro & segundo parâmetro & $1/a$ & $r_0$ \\
\midrule
\TabEtrychUnitarioCorpo
\bottomrule
\end{tabular}

\vspace{0.45cm}

\begin{tabular}{rrlrrrrr}
\toprule
$r_\mathrm{eff}$ & $a_-$ & forma & primeiro & segundo & $a$ & $r_0$ & estados \\
\midrule
\TabEtrychAminusCorpo
\bottomrule
\end{tabular}
\end{table}

\FloatBarrier
\subsection{A feição do lado positivo}
\label{sec:aplus}

A Tabela~\ref{tab:iii} publica também $a_p$, a feição do lado $a>0$, para
quatro das cinco linhas. Ela foi calibrada pela mesma prescrição, e a razão para
fazê-lo é experimental: no lado positivo a feição de perda é a de átomo-dímero,
e é reportada como a mais acessível e robusta das duas medidas. $a_-$ fica no
lado de três átomos, onde o trímero cruza o contínuo.

Do ponto de vista do código não há nada novo --- $r_0/a$ é o único número de que
uma solução depende, e nada na máquina sabe de que lado da ressonância está. O
que muda é a física em que ela aterra: $a>0$ carrega um dímero raso, de tamanho
$\sim a$, que aparece na coluna de estados.

\begin{table}[htbp]
\centering
\caption{Calibração em $a=a_p$ com $r_0=r_\mathrm{eff}$: \EtrychPositivasN{}
combinações, em raios de Bohr. O pior desvio é \EtrychPositivaA{} em $a$ e
\EtrychPositivaRzero{} em $r_0$ --- a mesma qualidade do lado negativo. A linha
de $162{,}36$~G não publica $a_p$ e nada foi inventado para ela.}
\label{tab:positivas}
\small
\begin{tabular}{rrlrrrrr}
\toprule
$B_\mathrm{res}$ (G) & $a_p$ & forma & primeiro & segundo & $a$ & $r_0$ & estados \\
\midrule
\TabEtrychPositivasCorpo
\bottomrule
\end{tabular}
\end{table}

\paragraph{E vale a pena?} Só se o lado teórico não for pior condicionado lá,
porque uma medida melhor entregue a uma calibração pior não adianta nada. Não é:
a dependência de forma nos quatro pontos de $a_p$ fica entre
\EtrychLequePositivoBaixo{} e \EtrychLequePositivoAlto{}, contra
\EtrychLequeRaso{} a \EtrychLequeAmenos{} nos pontos de $a_-$. O par simétrico
fecha a questão --- \EtrychLequeAmenos{} em $r_\mathrm{eff}/a=-0{,}16$ contra
\EtrychLequeMaisDezesseis{} em $+0{,}16$: \textbf{a dependência de forma depende
de $|r_\mathrm{eff}/a|$, não do sinal.} O lado do $a_p$ parece marginalmente pior
apenas porque a linha de $472{,}33$~G tem $r_\mathrm{eff}/a_p$ um pouco maior em
módulo que as demais.

A conclusão prática é que, se o experimento mede $a_p$ melhor, o lado positivo é
a comparação melhor nas duas pontas, e não um substituto de segunda. Uma
ressalva de alcance experimental fica em aberto: qual das ressonâncias da
Tabela~\ref{tab:iii} cabe no campo máximo do aparato é dado do laboratório e não
foi confirmado aqui, de modo que nenhum teto de campo é afirmado neste
relatório.

\begin{table}[htbp]
\centering
\caption{\textbf{Item (ii).} O que cada prescrição faz com $r_0/r_\mathrm{eff}$ ao
longo da mesma caminhada, de $r_\mathrm{eff}/a=\EtrychAlvoBaixo$ a
$+\EtrychAlvoAlto$, lida nos mesmos valores de $1/a$ para todas as formas.}
\label{tab:fam}
\small
\begin{tabular}{lrrrrrr}
\toprule
forma & \multicolumn{3}{c}{só a intensidade} & \multicolumn{2}{c}{ambos reajustados} & intensidade \\
\cmidrule(lr){2-4}\cmidrule(lr){5-6}
 & menor $r_0$ & maior $r_0$ & abertura & menor $r_0$ & maior $r_0$ & necessária \\
\midrule
\TabEtrychFamiliasCorpo
\bottomrule
\end{tabular}
\end{table}

Três leituras da Tabela~\ref{tab:fam} e das figuras:

\begin{enumerate}\itemsep2pt
  \item \textbf{Duas dispersões distintas, que não devem ser confundidas.} A
        coluna ``abertura'' é quanto o $r_0$ de \emph{uma mesma} forma anda ao
        longo de toda a varredura --- de \EtrychDerivaWell{} a
        \EtrychDerivaLj{} nas quatro formas, e \EtrychDerivaCseis{} na
        caminhada por $C_6$ --- e é propriedade da prescrição 2, que, ao
        congelar o alcance, deixa de manter $r_0$ fixo. A \emph{dependência de
        forma} é outra coisa: a discordância \emph{entre} as quatro formas no
        \emph{mesmo} ponto. Ela vale zero na unitariedade por construção,
        cresce quase linearmente a partir dali e atinge apenas
        \EtrychLequeAmenos{} em $r_\mathrm{eff}/a=-0{,}16$, extremo da faixa onde
        os $a_-$ se encontram, e \EtrychLequeRaso{} em $-0{,}06$. É essa a
        incerteza que um cálculo de três corpos herdaria da escolha de forma,
        e ela é pequena comparada com o fator $1{,}3$ a $1{,}4$ que separa o
        $a_-$ medido ($-12{,}1$ a $-14{,}1\,r_\mathrm{vdW}$) da previsão de van
        der Waals ($-9{,}7$). Resta medir o fator com que ela se propaga de
        $r_0$ para o $a_-$ previsto, que não é unitário e não foi determinado
        aqui.
  \item \textbf{O Lennard-Jones anda menos.} Precisa de \EtrychForcaLj{} de
        variação na intensidade contra \EtrychForcaWell{} do poço quadrado
        para a mesma faixa de $1/a$; por $C_6$ sozinho, apenas
        \EtrychForcaCseis{}, ao preço de o mínimo do potencial deslizar
        \EtrychMinimoCseis{} ao longo do caminho em vez de ficar parado.
  \item \textbf{Os estados ligados} contam $0$ em todo o lado $a<0$ e $1$ logo
        após a ressonância, nas quatro formas e nas duas prescrições --- o
        dímero raso entrando. Nenhuma sequência pula de ramo, verificado ponto
        a ponto.
\end{enumerate}

\begin{figure}[htbp]
\centering
\includegraphics[width=0.70\textwidth]{../../resultados/figuras/alcance_efetivo.pdf}\\[2pt]
\includegraphics[width=0.70\textwidth]{../../resultados/figuras/intensidade.pdf}\\[2pt]
\includegraphics[width=0.70\textwidth]{../../resultados/figuras/estados_ligados.pdf}
\caption{Acima: alcance efetivo ao longo do caminho --- tracejado, prescrição 1
(com $r_0$ fixo por construção); cheio, prescrição 2; pontilhado roxo, a
caminhada com $C_6$ sozinho. Meio: o parâmetro de intensidade, em unidades do
valor unitário. Abaixo: a contagem de estados ligados, com larguras de traço
escalonadas porque as quatro formas coincidem; o degrau aparece um ponto de
malha após a origem porque o ponto em $1/a=0$ é a semente unitária, onde o
estado está exatamente no limiar --- é resolução da varredura, não física.
Verticais pontilhadas: cinzas os $a_-$, amarelas os $a_p$.}
\label{fig:tudo}
\end{figure}

\section{A premissa, medida}
\label{sec:premissa}

Todo o resto deste relatório depende de uma única afirmação: que em energia
baixa a colisão guarda $a$ e $r_0$ e mais nada do potencial. É por isso que
calibrar duas formas diferentes nos mesmos dois números produz a mesma física,
e é por isso que entregar $(a_-,r_\mathrm{eff})$ a um cálculo de três corpos faz
sentido. A afirmação é padrão --- e, até aqui, era \emph{suposta}. Todo número
deste trabalho é calculado em $E=0$, onde $a$ e $r_0$ são \emph{definidos};
a premissa, porém, é uma afirmação sobre energia finita. Suposição e verificação
não são a mesma coisa, então ela foi medida.

\paragraph{Como.} As quatro formas são calibradas no ponto de $a_-$, de modo a
compartilhar $a$ e $r_0$ exatamente, e então cada uma é integrada em
$E=k^2/2$ --- o fator vem de $u''=2(V-E)u$, cujas soluções livres vão como
$\sin(kr+\delta)$ com $k^2=2E$. Casando com essa forma no corte sai
$k\cot\delta$, que a expansão de alcance efetivo prevê ser
$-1/a+r_0k^2/2$. Cada forma carrega um deslocamento constante próprio, de
origem numérica --- a fronteira interna está em raio finito e aqui o corte não é
extrapolado --- e esse deslocamento não é física. Ele é removido tomando a
diferença contra um $k$ de referência pequeno, e o termo de alcance efetivo é
subtraído. \textbf{O que sobra é exatamente a colisão que $a$ e $r_0$ não
determinam.}

\begin{table}[htbp]
\centering
\caption{O que $(a,r_0)$ \emph{não} fixa: espalhamento de $k\cot\delta$ entre
as quatro formas calibradas no mesmo $(a,r_0)$, no ponto de $a_-$. Comparar com
o próprio $k\cot\delta$, que ali vale $-1/a=\EtrychAlvoRazaoMod{}$. Em unidades
de $r_0=1$.}
\label{tab:premissa}
\small
\begin{tabular}{cc}
\toprule
$k\,r_0$ & espalhamento entre as quatro formas \\
\midrule
\TabEtrychPremissaCorpo
\bottomrule
\end{tabular}
\end{table}

\paragraph{O resultado.} Quatro potenciais tão diferentes quanto um degrau
quadrado, um poço $\mathrm{sech}^2$ suave, uma gaussiana e um Lennard-Jones com
cauda de van der Waals --- casados apenas em dois números --- concordam no
deslocamento de fase a \EtrychPremissaBaixo{} em
$k r_0=\EtrychPremissaKbaixo{}$ e a \EtrychPremissaMeio{} em
$k r_0=\EtrychPremissaKmeio{}$. A premissa se sustenta.

E sustenta-se de um jeito informativo, não apenas aprovado. O que sobra
\emph{não} é zero, e não deveria ser: a expansão continua, e o termo seguinte é
o parâmetro de forma em $k^4$. Medido, o resíduo segue potência
\textbf{\EtrychPremissaPotencia{}} com coeficiente \EtrychPremissaCoef{}. Uma
potência de exatamente quatro é a assinatura do termo físico; um deslocamento
de casamento mal feito não encolheria com $k$ de modo algum. Ou seja: não só a
premissa vale, como o primeiro termo que ela ignora está identificado e
quantificado.

\paragraph{E no setor ligado.} A premissa costuma ser enunciada sobre colisões,
mas o dímero é a outra metade dela. Do lado $a>0$ existe um estado ligado raso,
e as quatro formas calibradas em $r_0/a=\EtrychDimeroRazao{}$ o põem na mesma
energia a \EtrychDimeroLeque{} --- $\EtrychDimeroEnergia{}$ em unidades de
$r_0=1$, igual em cinco casas nas quatro. Os estados ligados são contados aqui
por um caminho que não olha $E=0$ em momento algum: varre-se $E<0$ procurando
zeros do resíduo de casamento do estado ligado, o que não compartilha nada com
a contagem de nós além do integrador. As duas contagens concordam nas quatro
formas e em todos os ramos.

\paragraph{Por que isso é a conclusão central.} A dependência de forma de
\EtrychLequeAmenos{} medida na Seção~\ref{sec:resultados} diz quanto as formas
discordam \emph{no} $r_0$. Esta seção diz algo mais forte: fixados $a$ e $r_0$,
não sobra praticamente nada para discordar. As duas juntas é que fecham o
argumento de que entregar $(a_-,r_\mathrm{eff})$ a três corpos é um procedimento
bem posto --- a primeira limita o erro de entrada, a segunda mostra que não há
um canal escondido pelo qual a forma volte a importar.

\section{O ramo: quantos dímeros profundos o potencial carrega}
\label{sec:ramos}

Um resultado que não estava no escopo da tarefa e que organiza o passo
seguinte. \textbf{O par $(a,r_0)$ não determina o potencial.} Aprofundando uma
forma além de cada polo de $a$, um novo estado ligado aparece e o mesmo
$(a,r_0)$ é alcançado de novo, agora com um dímero profundo a mais. A
calibração tem, portanto, uma solução por ramo.

\begin{table}[htbp]
\centering
\caption{A mesma calibração --- $a_-=-830$, $r_\mathrm{eff}=136\,a_0$ --- em três
ramos do poço quadrado. Os três são indistinguíveis em espalhamento de dois
corpos de baixa energia e carregam zero, um e dois dímeros profundos. No poço
quadrado o alcance quase não muda --- \EtrychRamoPocoMove{} entre os três
ramos --- e o que muda é a profundidade. Isso é propriedade \emph{do poço}, e
não do fenômeno: medido nas outras formas, o alcance varia por um fator de
até \EtrychRamoOutrasMove{} entre os mesmos três ramos.}
\label{tab:ramos}
\small
\begin{tabular}{crrrrr}
\toprule
dímeros profundos & $v$ & $R$ ($a_0$) & $R/r_\mathrm{vdW}$ & $a$ ($a_0$) & $r_0$ ($a_0$) \\
\midrule
\TabEtrychRamosCorpo
\bottomrule
\end{tabular}
\end{table}

\paragraph{Por que isso importa, e por que é testável.} Em dois corpos os três
potenciais da Tabela~\ref{tab:ramos} são indistinguíveis --- é exatamente o que
a universalidade afirma. Em três corpos não são, por uma razão concreta: um
dímero profundo dá ao trímero de Efimov \emph{para onde decair}. A taxa desse
decaimento é o que o parâmetro de largura $\eta$ mede. No ramo mais raso não há
dímeros profundos, não há esse canal, e $\eta$ seria exatamente zero.

A Tabela~\ref{tab:iii} publica $\eta_-$ para todas as cinco linhas, entre
\EtrychEtaMenor{} e \EtrychEtaMaior{}, todos claramente diferentes de zero ---
como se espera do potássio-39, que carrega dezenas de estados moleculares
profundos. \textbf{A medida já exclui o ramo raso como descrição do sistema
real.}

Isso converte a escolha do ramo, que no escopo original era uma decisão
arbitrária registrada como pergunta em aberto, em um \emph{teste}: um cálculo
de três corpos em cada ramo pode ser comparado com o $\eta_-$ publicado. O ramo
raso passa a ser o caso de controle --- o que a universalidade de alcance zero
sozinha diria --- e um ramo com estados profundos é o candidato realista. Se um
deles reproduzir $\eta_-\approx0{,}25$ a $0{,}3$, isso valida o potencial por
um caminho independente do próprio $a_-$.

\paragraph{Como o ramo é localizado.} A fronteira entre ramos é onde um novo
estado liga, e há duas maneiras de perguntar isso que \emph{não} concordam. A
contagem de estados ligados é uma escada não decrescente na profundidade ---
aprofundar um potencial não desliga estado --- então cada degrau é o único
cruzamento do teste ``contagem $\geq n$'', que a bisseção localiza e não pode
pular, por estreita que a janela seja. Mas o degrau não é a fronteira física: a
contagem só soma o nó que a reta assintótica contribui em $r=a$ depois que
$|r_0/a|$ passa de um limiar, e portanto atrasa. Medido, o atraso é
\EtrychEscadaAtraso{} na profundidade, o mesmo nas quatro fronteiras do poço
quadrado --- artefato sistemático de limiar, não ruído.

A fronteira física é o \emph{polo de $a$}, onde $1/a=0$ e o novo estado chega
com energia de ligação exatamente nula. Nenhum limiar entra nela. Então o degrau
serve só para encurralar, e a fronteira é o zero de $1/a$ dentro desse cerco.
Isso é conferível, e não apenas convergente: para duas das quatro formas a
escada de polos é analítica, $\EtrychEscadaPoco{}$ no poço quadrado e
$\EtrychEscadaMpt{}$ no Pöschl--Teller, e a bisseção as reproduz a
\EtrychEscadaErro{} relativo sem nada ajustado.

Dentro de uma janela, $1/a$ corre monotonicamente de $0^+$ na borda de baixo a
$0^-$ na de cima, o que só consegue fazer passando por infinito: $a$ cruza zero
exatamente uma vez e $r_0$ diverge com ele, de modo que $r_0/a$ tem um polo
dentro de cada janela. A janela é partida ali, e cada metade é varrida até a
contagem de raízes repetir --- $r_0/a$ não é monótona ali, sobe, vira e cruza
zero, e um alvo pode ser atingido mais de uma vez.

\paragraph{Dois detalhes que a implementação expôs.} O primeiro: o ramo não
determina o potencial sozinho. A busca encontra mais de uma solução dentro do
mesmo ramo, com o mesmo $(a,r_0)$ e o mesmo número de estados, diferindo no
alcance --- no ramo de um dímero profundo, o poço quadrado admite $136\,a_0$ e
também $1428\,a_0$. A segunda é descartada por critério físico, não numérico:
$1428\,a_0$ são vinte e dois alcances de van der Waals, e nada no potássio-39
é tão longo. As soluções mantidas ficam entre \EtrychRamoAlcanceBaixo{} e
\EtrychRamoAlcanceAlto{} vezes $r_\mathrm{vdW}$.

O segundo: exigir $r_0>0$ não é consequência, é requisito. A razão $r_0/a$
atinge o alvo também com os dois sinais invertidos, e reescalar por $1/r_0$
pediria então um fator de escala negativo. A exigência física é
$r_0=+r_\mathrm{eff}$ com $a=a_-<0$, e reescalar por fator positivo não conserta
sinal.

\section{Verificação}

Quarenta e sete conferências automáticas, agrupadas por \emph{contra o que}
comparam --- distinção que importa, porque conferir contra forma fechada não é
conferir o integrador, e conferir o código contra ele mesmo não é conferir
correção.

\begin{table}[htbp]
\centering
\caption{As conferências de \cod{test\_lab.py}, por natureza da âncora.}
\small
\begin{tabular}{lc@{\quad}p{7.1cm}}
\toprule
grupo & n & âncora \\
\midrule
formas fechadas & 9 & $a$, $r_0$ e estado ligado do poço, poço unitário,
  família, as duas escadas de polos \\
tabelas publicadas & 6 & hélio com quatro potenciais, contagem de nós, mínimos
  de Aziz, constantes \\
segundo método & 3 & o Numerov, e a contagem de autovalores contra a de nós \\
invariância e estrutura & 9 & escala, corte do caroço, fim da malha, sinal e
  limiar na unitariedade, simetria do leque \\
a premissa & 4 & deslocamento de fase, as quatro formas em energia finita, a
  energia do dímero \\
a própria entrega & 14 & alvos nos dois lados, cauda, ramos, recusas \\
o código contra si mesmo & 2 & convergência de malha, pino de regressão \\
\bottomrule
\end{tabular}
\end{table}

Contra as formas fechadas do poço quadrado o desvio relativo é \EtrychPocoA{}
em $a$ e \EtrychPocoRzero{} em $r_0$. Nas vinte linhas da calibração, o pior
desvio dos alvos é \EtrychAlvoA{} em $a$ e \EtrychAlvoRzero{} em $r_0$; nas
quatro soluções unitárias, $1/a$ fica em \EtrychUnitInv{} e $r_0$ dista
\EtrychUnitRzero{} de $1$.

O Numerov em malha logarítmica, que era o caminho de produção na versão
anterior, virou conferência independente: ordem $4$ contra $8$, passo fixo em
$\ln r$ contra adaptativo em $r$, Simpson em malha própria contra integral
dentro do sistema. A troca não foi por precisão de integração --- nela os dois
empatam --- e sim pela Eq.~\eqref{eq:ar0}: por quadratura sobre a solução,
aquela diferença de números quase iguais sofre cancelamento, e refinar a malha
\emph{piora} a resposta ($7\times10^{-6}$ com $8\,001$ pontos,
$5{,}6\times10^{-3}$ com $128\,001$). Um método que piora quando refinado não
consegue informar o próprio erro.

O maior erro \emph{numérico} que resta não é mais o integrador, e sim a cauda
de van der Waals de segunda ordem do Lennard-Jones, que a extrapolação remove
apenas em primeira ordem; as outras três formas já concordam nos três cortes.
O tamanho dele foi medido: extrapolar de $(R,3R)$ e de $(3R,9R)$ --- que não
compartilham corte algum --- discorda em \EtrychCaudaResiduo{} relativo. Isso
não é o que limita o resultado. A dependência de forma no ponto de $a_-$ é
\EtrychLequeAmenos{}, mais de quatro ordens de grandeza acima, e é \emph{ela} o
limitante. Por isso nenhuma correção analítica de cauda foi implementada:
trocaria \EtrychCaudaResiduo{} por algo menor num cálculo cujo sinal é
\EtrychLequeAmenos{}. Empurrar o corte para mais longe também não ajuda, e
piora: de $(9R,27R)$ a discordância sobe para \EtrychCaudaResiduoLonge{},
porque a reta assintótica cresce como $r$ e carrega o arredondamento
acumulado.

\section{Limitações}

\textbf{Canal único.} Um potencial de canal único não reproduz integralmente
uma ressonância de Feshbach, que é problema de dois canais. As prescrições
imitam a variação de $a$, não o acoplamento.

\textbf{Nenhuma das ressonâncias é larga.} A ressalva pesa mais do que
pareceria, porque a universalidade de van der Waals para o parâmetro de três
corpos é derivada no regime $s_\mathrm{res}\gg1$, e nenhuma das cinco linhas está
nele --- os autores classificam quatro delas como de intensidade intermediária.
Sinal independente disso: eles próprios observam que, para as quatro mais
fortes, o $r_\mathrm{eff}$ fica cerca de $25\%$ abaixo do
$\approx2{,}8\,r_\mathrm{vdW}$ esperado para ressonâncias largas. O alcance
efetivo que entra nesta calibração já não é o de ressonância larga --- o que
torna a correção de alcance finito ainda mais pertinente, e é exatamente ela
que este trabalho embute no potencial.

\textbf{A linha de $162{,}36$~G} é a que mais destoa em tudo: $s_\mathrm{res}$ e
$r_\mathrm{eff}$ bem menores que os demais e $r_\mathrm{eff}/a_-$ cerca de metade. É
ali que o caráter de dois canais deve aparecer primeiro.

\textbf{O que não foi medido.} O fator com que a incerteza de forma se propaga
de $r_0$ para o $a_-$ previsto. Não é unitário, e só o cálculo de três corpos
o determina.

\paragraph{Reprodução.} \cod{python test\_lab.py} roda as conferências;
\cod{python notas\_teoria/fonte/gerar.py} regenera as tabelas, os números
citados no texto, as três figuras e os arquivos \cod{calibracoes.csv},
\cod{familias.csv} e \cod{corte\_da\_cauda.csv} em \cod{resultados/}.

\begin{thebibliography}{9}\small
\bibitem{etrych}
J.~Etrych, G.~Martirosyan, A.~Cao, J.~A.~P.~Glidden, L.~H.~Dogra,
J.~M.~Hutson, Z.~Hadzibabic e C.~Eigen,
\textit{Pinpointing Feshbach resonances and testing Efimov universalities in
$^{39}$K}, Phys.\ Rev.\ Research \textbf{5}, 013174 (2023).

\bibitem{macedo}
A.~Macêdo-Lima e L.~Madeira,
\textit{Scattering length and effective range of low-energy collisions},
Rev.\ Bras.\ Ensino Fís.\ \textbf{45}, e20230079 (2023).

\bibitem{motovilov}
A.~K.~Motovilov, W.~Sandhas, S.~A.~Sofianos e E.~A.~Kolganova,
\textit{Binding energies and scattering observables in the $^4$He$_3$ atomic
system}, Eur.\ Phys.\ J.\ D \textbf{13}, 33 (2001).
\end{thebibliography}

\end{document}
